\chapter{1994年全国卷（文）}
\section{单选题}

\begin{problem}
设全集 \(I = \{0, 1, 2, 3, 4\}\)，集合 \(A = \{0, 1, 2, 3\}\)，集合 \(B = \{2, 3, 4\}\)，则 \(\overline{A} \cup \overline{B} =\)（\quad）

\choices
    {\(\{0\} \)}
    {\(\{0,1\}\)}
    {\(\{0,1,4\}\)}
    {\(\{0,1,2,3,4\}\)}
\end{problem}

\begin{answer}
C
\end{answer}

\begin{solution}
补集均相对于全集 \(I\) 而言，因此 \(\overline{A}=\{4\}\)，\(\overline{B}=\{0,1\}\). 所以 \(\overline{A}\cup\overline{B}=\{0,1,4\}\)，故选 C.
\end{solution}

\begin{problem}
如果方程 \(x^{2} + ky^{2} = 2\) 表示焦点在 \(y\) 轴上的椭圆，那么实数 \(k\) 的取值范围是（\quad）

\choices
    {\((0, +\infty)\)}
    {\((0,2)\)}
    {\((1, +\infty)\)}
    {\((0,1)\)}
\end{problem}

\begin{answer}
D
\end{answer}

\begin{solution}
要表示椭圆，必须有 \(k>0\). 方程化为 \(\dfrac{x^{2}}{2}+\dfrac{y^{2}}{2/k}=1\)，其两个半轴长的平方分别为 \(2\) 和 \(2/k\). 焦点在 \(y\) 轴上要求 \(y\) 轴方向的半轴较长，即 \(2/k>2\)，从而 \(0<k<1\). 故选 D.
\end{solution}

\begin{problem}
点 \((0,5)\) 到直线 \(y=2x\) 的距离是（\quad）

\choices
    {\(\frac{5}{2}\)}
    {\(\sqrt{5}\)}
    {\(\frac{3}{2}\)}
    {\(\frac{\sqrt{5}}{2}\)}
\end{problem}

\begin{answer}
B
\end{answer}

\begin{solution}
直线写成 \(2x-y=0\). 由点到直线的距离公式，距离为 \(\dfrac{|2\cdot0-5|}{\sqrt{2^{2}+(-1)^{2}}}=\dfrac{5}{\sqrt5}=\sqrt5\)，故选 B.
\end{solution}

\begin{problem}
设 \(\theta\) 是第二象限的角，则必有（\quad）

\choices
    {\(\tan \frac{\theta}{2} > \cot \frac{\theta}{2}\)}
    {\(\tan \frac{\theta}{2} < \cot \frac{\theta}{2}\)}
    {\(\sin \frac{\theta}{2} >\cos \frac{\theta}{2}\)}
    {\(\sin \frac{\theta}{2} < \cos \frac{\theta}{2}\)}
\end{problem}

\begin{answer}
A
\end{answer}

\begin{solution}
设 \(\theta=2k\pi+\alpha\)，其中 \(k\in\Z\)，且 \(\dfrac{\pi}{2}<\alpha<\pi\). 则 \(\dfrac{\theta}{2}=k\pi+\dfrac{\alpha}{2}\)，其中 \(\dfrac{\pi}{4}<\dfrac{\alpha}{2}<\dfrac{\pi}{2}\). 因为正切函数的周期为 \(\pi\)，所以 \(\tan\dfrac{\theta}{2}=\tan\dfrac{\alpha}{2}>1\)，从而 \(\cot\dfrac{\theta}{2}<1\)，必有 \(\tan\dfrac{\theta}{2}>\cot\dfrac{\theta}{2}\)，答案为 A. 当 \(k\) 为偶数时，\(\dfrac{\theta}{2}\) 在第一象限，\(\sin\dfrac{\theta}{2}>\cos\dfrac{\theta}{2}\)；当 \(k\) 为奇数时，\(\dfrac{\theta}{2}\) 在第三象限，\(\sin\dfrac{\theta}{2}<\cos\dfrac{\theta}{2}\)，所以后两个选项均非必然成立.
\end{solution}

\begin{problem}
某种细菌在培养过程中，每\(20\text{ 分钟}\)分裂一次（\(1\text{ 个}\)分裂为\(2\text{ 个}\)）. 经过\(3\text{ 小时}\)，这种细菌由\(1\text{ 个}\)可繁殖成（\quad）

\choices
    {\(511\text{ 个}\)}
    {\(512\text{ 个}\)}
    {\(1023\text{ 个}\)}
    {\(1024\text{ 个}\)}
\end{problem}

\begin{answer}
B
\end{answer}

\begin{solution}
\(3\text{ 小时}=180\text{ 分钟}\)，共分裂 \(180\div20=9\) 次. 每次分裂后总数乘以 \(2\)，所以最终个数为 \(1\cdot2^{9}=512\)，故选 B.
\end{solution}

\begin{problem}
在下列函数中，以 \(\frac{\pi}{2}\) 为周期的函数是（\quad）

\choices
    {\(y = \sin 2x + \cos 4x\)}
    {\(y = \sin 2x\cos 4x\)}
    {\(y = \sin 2x + \cos 2x\)}
    {\(y = \sin 2x\cos 2x\)}
\end{problem}

\begin{answer}
D
\end{answer}

\begin{solution}
将 \(x\) 换成 \(x+\dfrac{\pi}{2}\)，有 \(\sin(2x+\pi)=-\sin2x\)，\(\cos(4x+2\pi)=\cos4x\)，因此 A、B 均不保持原函数值；C 中两项均变号，也不保持原函数值. D 可化为 \(y=\dfrac12\sin4x\)，而 \(\dfrac12\sin\left(4x+2\pi\right)=\dfrac12\sin4x\)，故 \(\dfrac{\pi}{2}\) 是其周期，选 D.
\end{solution}

\begin{problem}
已知正六棱台的上、下底面边长分别为 \(2\) 和 \(4\)，高为 \(2\)，则其体积为（\quad）

\choices
    {\(32\sqrt{3}\)}
    {\(28\sqrt{3}\)}
    {\(24\sqrt{3}\)}
    {\(20\sqrt{3}\)}
\end{problem}

\begin{answer}
B
\end{answer}

\begin{solution}
边长为 \(s\) 的正六边形面积为 \(\dfrac{3\sqrt3}{2}s^{2}\)，故上、下底面积分别为 \(S_1=6\sqrt3\) 和 \(S_2=24\sqrt3\). 棱台体积为
\(V=\dfrac{h}{3}\left(S_1+S_2+\sqrt{S_1S_2}\right)=\dfrac{2}{3}\left(6\sqrt3+24\sqrt3+12\sqrt3\right)=28\sqrt3\). 故选 B.
\end{solution}

\begin{problem}
设 \(F_{1}\) 和 \(F_{2}\) 为双曲线 \(\frac{x^{2}}{4} - y^{2} = 1\) 的两个焦点，点 \(P\) 在双曲线上且满足 \(\angle F_{1}PF_{2} = 90^{\circ}\)，则 \(\triangle F_{1}F_{2}P\) 的面积是（\quad）

\choices
    {\(1\)}
    {\(\frac{\sqrt{5}}{2}\)}
    {\(2\)}
    {\(\sqrt{5}\)}
\end{problem}

\begin{answer}
A
\end{answer}

\begin{solution}
双曲线的半实轴、半虚轴分别为 \(a=2\)、\(b=1\)，所以焦半距 \(c=\sqrt{a^{2}+b^{2}}=\sqrt5\)，且 \(F_1F_2=2\sqrt5\). 令 \(u=PF_1\)、\(v=PF_2\). 由双曲线定义，\(|u-v|=2a=4\)；又因 \(\angle F_1PF_2=90^{\circ}\)，由勾股定理有 \(u^{2}+v^{2}=F_1F_2^{2}=20\). 因此 \(16=(u-v)^2=u^2+v^2-2uv=20-2uv\)，得 \(uv=2\). 三角形面积为 \(\dfrac12uv=1\)，故选 A.
\end{solution}

\begin{problem}
如果复数 \(z\) 满足 \(|z + \i| + |z - \i| = 2\)，那么 \(|z + \i + 1|\) 的最小值是（\quad）

\choices
    {\(1\)}
    {\(\sqrt{2}\)}
    {\(2\)}
    {\(\sqrt{5}\)}
\end{problem}

\begin{answer}
A
\end{answer}

\begin{solution}
在复平面中，\(-\i\) 和 \(\i\) 的距离为 \(2\). 由三角不等式，等式 \(|z+\i|+|z-\i|=2\) 成立当且仅当 \(z\) 在线段 \([-\i,\i]\) 上，因此可设 \(z=t\i\)，其中 \(-1\leqslant t\leqslant1\). 于是
\(|z+\i+1|=|1+(t+1)\i|=\sqrt{1+(t+1)^2}\geqslant1\)，且在 \(t=-1\) 时取到，故选 A.
\end{solution}

\begin{problem}
有甲、乙、丙三项任务，甲需\(2\) 人承担，乙、丙各需\(1\) 人承担. 从\(10\) 人中选派\(4\) 人承担这三项任务，不同的选法共有（\quad）

\choices
    {\(1260\text{ 种}\)}
    {\(2025\text{ 种}\)}
    {\(2520\text{ 种}\)}
    {\(5040\text{ 种}\)}
\end{problem}

\begin{answer}
C
\end{answer}

\begin{solution}
甲、乙、丙三项任务有区别. 先从 \(10\) 人中选出承担甲的 \(2\) 人，再从余下 \(8\) 人中选承担乙的 \(1\) 人，最后从余下 \(7\) 人中选承担丙的 \(1\) 人，选法数为 \(\mathrm C_{10}^{2}\cdot8\cdot7=45\cdot8\cdot7=2520\)，故选 C.
\end{solution}

\begin{problem}
对于直线 \(m\)，\(n\) 和平面 \(\alpha\)，\(\beta\)，\(\alpha \perp \beta\) 的一个充分条件是（\quad）

\choices
    {\(m \perp n\)，\(m \parallel \alpha\)，\(n \parallel \beta\)}
    {\(m \perp n\)，\(\alpha \cap \beta = m\)，\(n \subset \alpha\)}
    {\(m \parallel n\)，\(n \perp \beta\)，\(m \subset \alpha\)}
    {\(m \parallel n\)，\(m \perp \alpha\)，\(n \perp \beta\)}
\end{problem}

\begin{answer}
C
\end{answer}

\begin{solution}
若选项 C 成立，则 \(m\parallel n\) 且 \(n\perp\beta\)，所以 \(m\perp\beta\). 又因为 \(m\subset\alpha\)，平面 \(\alpha\) 含有一条垂直于平面 \(\beta\) 的直线，故 \(\alpha\perp\beta\). 因此选项 C 是充分条件，故选 C.
\end{solution}

\begin{problem}
设函数 \(f(x) = 1 - \sqrt{1 - x^2}\)（\(-1 \leqslant x \leqslant 0\)），则函数 \(y = f^{-1}(x)\) 的图象是（\quad）

\choices
    {\choicebitmap[width=0.15\paperwidth]{img/1994/national_liberal/q12_optA.png}}
    {\choicebitmap[width=0.15\paperwidth]{img/1994/national_liberal/q12_optB.png}}
    {\choicebitmap[width=0.15\paperwidth]{img/1994/national_liberal/q12_optC.png}}
    {\choicebitmap[width=0.15\paperwidth]{img/1994/national_liberal/q12_optD.png}}
\end{problem}

\begin{answer}
B
\end{answer}

\begin{solution}
令 \(u=f(t)=1-\sqrt{1-t^2}\)，其中 \(-1\leqslant t\leqslant0\). 则 \(0\leqslant u\leqslant1\)，且 \(\sqrt{1-t^2}=1-u\)，从而 \(t^2=2u-u^2\). 由于 \(t\leqslant0\)，有 \(t=-\sqrt{2u-u^2}\). 所以反函数图象满足 \(y=-\sqrt{2x-x^2}\)，即 \((x-1)^2+y^2=1\)，其中 \(0\leqslant x\leqslant1\)、\(-1\leqslant y\leqslant0\)，是位于第四象限的四分之一圆弧，故选 B.
\end{solution}

\begin{problem}
已知过球面上 \(A\)，\(B\)，\(C\) 三点的截面和球心的距离等于球半径的一半，且 \(AB = BC = CA = 2\)，则球面面积是（\quad）

\choices
    {\(\frac{16}{9}\pi\)}
    {\(\frac{8}{3}\pi\)}
    {\(4\pi\)}
    {\(\frac{64}{9}\pi\)}
\end{problem}

\begin{answer}
D
\end{answer}

\begin{solution}
截面内的三点构成边长为 \(2\) 的正三角形，其外接圆半径为 \(r=\dfrac{2}{\sqrt3}\). 设球半径为 \(R\)，截面与球心的距离为 \(d=\dfrac R2\). 截面圆半径满足 \(r^2=R^2-d^2=\dfrac34R^2\)，所以 \(R=\dfrac{2r}{\sqrt3}=\dfrac43\). 球面面积为 \(4\pi R^2=\dfrac{64}{9}\pi\)，故选 D.
\end{solution}

\begin{problem}
如果函数 \(y = \sin 2x + a \cos 2x \) 的图象关于直线 \(x = -\frac{\pi}{8} \) 对称，那么 \(a =\)（\quad）

\choices
    {\(\sqrt{2}\)}
    {\(-\sqrt{2}\)}
    {\(1\)}
    {\(-1\)}
\end{problem}

\begin{answer}
D
\end{answer}

\begin{solution}
设 \(y=R\sin(2x+\varphi)\)，则 \(R>0\)，且 \(a=R\sin\varphi\)、\(1=R\cos\varphi\)，所以 \(\tan\varphi=a\). 正弦曲线的对称轴经过极值点，因此在对称轴 \(x=-\dfrac\pi8\) 上有 \(2x+\varphi=\dfrac\pi2+k\pi\)，即 \(\varphi=\dfrac{3\pi}{4}+k\pi\). 于是 \(a=\tan\varphi=-1\)，故选 D.
\end{solution}

\begin{problem}
定义在 \((- \infty, + \infty)\) 上的任意函数 \(f(x)\) 都可以表示成一个奇函数 \(g(x)\) 和一个偶函数 \(h(x)\) 之和，如果 \(f(x) = \lg (10^{x} + 1)\)，\(x \in (-\infty, +\infty)\)，那么（\quad）

\choices
    {\(g(x) = x\)，\(h(x) = \lg (10^{x} + 10^{-x} + 2)\)}
    {\(g(x) = \frac{1}{2}\bigl[\lg (10^{x} + 1) + x\bigr]\)，\(h(x) = \frac{1}{2}\bigl[\lg (10^{x} + 1) - x\bigr]\)}
    {\(g(x) = \frac{x}{2}\)，\(h(x) = \lg (10^x + 1) - \frac{x}{2}\)}
    {\(g(x) = -\frac{x}{2}\)，\(h(x) = \lg (10^x + 1) + \frac{x}{2}\)}
\end{problem}

\begin{answer}
C
\end{answer}

\begin{solution}
有 \(f(-x)=\lg(10^{-x}+1)=\lg\dfrac{10^x+1}{10^x}=f(x)-x\). 奇、偶分解唯一给出 \(g(x)=\dfrac{f(x)-f(-x)}2\)、\(h(x)=\dfrac{f(x)+f(-x)}2\). 代入上式，得 \(g(x)=\dfrac x2\)，\(h(x)=f(x)-\dfrac x2=\lg(10^x+1)-\dfrac x2\)，故选 C.
\end{solution}

\section{填空题}

\begin{problem}
在 \((3 - x)^7\) 的展开式中，\(x^5\) 的系数是\fillinblank{}.
\end{problem}

\begin{answer}
\(-189\)
\end{answer}

\begin{solution}
通项中取 \((-x)^5\) 的项才能得到 \(x^5\)，所以其系数为 \(\mathrm C_7^5\cdot3^{7-5}\cdot(-1)^5=21\cdot9\cdot(-1)=-189\).
\end{solution}

\begin{problem}
抛物线 \(y^{2} = 8 - 4x \) 的准线方程是\fillinblank{}，圆心在该抛物线的顶点且与其准线相切的圆的方程是\fillinblank{}.
\end{problem}

\begin{answer}
\(x=3\)，\(\left(x-2\right)^2+y^2=1\)
\end{answer}

\begin{solution}
方程化为 \(y^2=-4(x-2)\)，与标准式 \(y^2=4p(x-2)\) 比较得 \(p=-1\). 顶点为 \((2,0)\)，准线为 \(x=2-p=3\). 以顶点为圆心且与直线 \(x=3\) 相切的圆半径等于圆心到准线的距离 \(1\)，故圆方程为 \((x-2)^2+y^2=1\).
\end{solution}

\begin{problem}
已知 \(\sin \theta + \cos \theta = \frac{1}{5}\)，\(\theta \in (0, \pi)\)，则 \(\cot \theta\) 的值是\fillinblank{}.
\end{problem}

\begin{answer}
\(-\frac{3}{4}\)
\end{answer}

\begin{solution}
平方得 \(1+2\sin\theta\cos\theta=\dfrac1{25}\)，所以 \(\sin\theta\cos\theta=-\dfrac{12}{25}<0\). 又 \(\theta\in(0,\pi)\)，故 \(\sin\theta>0\)，从而 \(\cos\theta<0\). 令 \(s=\sin\theta\)、\(c=\cos\theta\)，则 \(s+c=\dfrac15\)、\(sc=-\dfrac{12}{25}\)，故 \(s\)，\(c\) 是方程 \(t^2-\dfrac15t-\dfrac{12}{25}=0\) 的两根，即 \(\{s,c\}=\left\{\dfrac45,-\dfrac35\right\}\). 由 \(s>0\) 得 \(s=\dfrac45\)、\(c=-\dfrac35\)，所以 \(\cot\theta=\dfrac c s=-\dfrac34\).
\end{solution}

\begin{problem}
设圆锥底面圆周上两点 \(A\)，\(B\) 间的距离为 \(2\)，圆锥顶点到直线 \(AB\) 的距离为 \(\sqrt{3}\)，\(AB\) 和圆锥的轴的距离为 \(1\)，则该圆锥的体积为\fillinblank{}.
\end{problem}

\begin{answer}
\(\frac{2\sqrt{2}}{3}\pi\)
\end{answer}

\begin{solution}
设底面圆心为 \(O\)，圆锥高为 \(h\)，底面半径为 \(r\). 轴与弦所在直线的距离等于 \(O\) 到 \(AB\) 的距离，因此在底面内，弦半长为 \(1\)，且 \(O\) 到弦的距离为 \(1\)，由直角三角形得 \(r^2=1^2+1^2=2\). 顶点到直线 \(AB\) 的垂距由轴向和底面内两个互相垂直的分量组成，所以 \(3=h^2+1^2\)，得 \(h=\sqrt2\). 于是圆锥体积为 \(V=\dfrac13\pi r^2h=\dfrac13\pi\cdot2\cdot\sqrt2=\dfrac{2\sqrt2}{3}\pi\).
\end{solution}

\begin{problem}
在测量某物理量的过程中，因仪器和观察的误差，使得 \(n\) 次测量分别得到 \(a_{1}\)，\(a_{2}\)，\(\cdots\)，\(a_{n}\)，共 \(n\) 个数据，我们规定所测量物理量的“最佳近似值” \(a\) 是这样一个量：与其他近似值比较，\(a\) 与各数据的差的平方和最小. 依此规定，从 \(a_{1}\)，\(a_{2}\)，\(\cdots\)，\(a_{n}\) 推出的 \(a =\fillinblank{}\).
\end{problem}

\begin{answer}
\(\frac{a_1+a_2+\cdots+a_n}{n}\)
\end{answer}

\begin{solution}
令 \(\bar a=\dfrac{a_1+a_2+\cdots+a_n}{n}\). 平方和可配方为
\(\sum_{i=1}^{n}(a-a_i)^2=\sum_{i=1}^{n}(\bar a-a_i)^2+n(a-\bar a)^2\). 第一项与 \(a\) 无关，第二项在且仅在 \(a=\bar a\) 时取最小值，因此最佳近似值为 \(a=\dfrac{a_1+a_2+\cdots+a_n}{n}\).
\end{solution}

\section{解答题}

\begin{problem}
求函数 \(y = \frac{\sin 3x\sin^3 x + \cos 3x\cos^3 x}{\cos^2 2x} + \sin 2x \) 的最小值.
\end{problem}

\begin{solution}
函数定义域要求 \(\cos2x\neq0\). 记 \(s=\sin x\)、\(c=\cos x\)，则分子
\(N=(3s-4s^3)s^3+(4c^3-3c)c^3=4(c^6-s^6)-3(c^4-s^4)\). 利用 \(s^2+c^2=1\)，有
\(N=(c^2-s^2)\left[4(c^4+c^2s^2+s^4)-3\right]=(c^2-s^2)(1-4s^2c^2)=\cos2x\cos^2 2x=\cos^3 2x\).

因此在定义域内，\(\dfrac{N}{\cos^2 2x}=\cos2x\)，原函数化为 \(y=\cos2x+\sin2x=\sqrt2\sin\left(2x+\dfrac\pi4\right)\). 故 \(y\geqslant-\sqrt2\). 当 \(x=\dfrac{5\pi}{8}\) 时，\(\cos2x\neq0\) 且 \(\sin\left(2x+\dfrac\pi4\right)=-1\)，等号可以取得，所以最小值为 \(-\sqrt2\).
\end{solution}

\begin{problem}
已知函数 \(f(x) = \log_a x\)（\(a > 0\) 且 \(a \neq 1\)，\(x \in \R_+\)），若 \(x_1\)，\(x_2 \in \R_+\)，判断 \(\frac{1}{2} [f(x_1) + f(x_2)] \) 与 \(f\left(\frac{x_1 + x_2}{2}\right) \) 的大小，并加以证明.
\end{problem}

\begin{solution}
由对数的运算性质，\(\dfrac12[f(x_1)+f(x_2)]=\dfrac12\log_a(x_1x_2)=\log_a\sqrt{x_1x_2}\). 又因为 \(x_1\)，\(x_2>0\)，由 \(\left(\sqrt{x_1}-\sqrt{x_2}\right)^2\geqslant0\) 得 \(2\sqrt{x_1x_2}\leqslant x_1+x_2\)，所以 \(\sqrt{x_1x_2}\leqslant\dfrac{x_1+x_2}{2}\)，且等号当且仅当 \(x_1=x_2\).

当 \(a>1\) 时，\(\log_a x\) 单调递增，所以 \(\dfrac12[f(x_1)+f(x_2)]\leqslant f\left(\dfrac{x_1+x_2}{2}\right)\). 当 \(0<a<1\) 时，\(\log_a x\) 单调递减，所以 \(\dfrac12[f(x_1)+f(x_2)]\geqslant f\left(\dfrac{x_1+x_2}{2}\right)\). 两种情形均在且仅在 \(x_1=x_2\) 时取等号.
\end{solution}

\begin{problem}
如图，已知 \(A_{1}B_{1}C_{1} - ABC\) 是正三棱柱，\(D\) 是 \(AC\) 中点.

\begin{enumerate}
\item 证明：\(AB_{1} \parallel\) 平面 \(DBC_{1}\)；
\item 假设 \(AB_{1} \perp BC_{1}\)，\(BC = 2\)，求线段 \(AB_{1}\) 在侧面 \(B_{1}BCC_{1}\) 上的射影长.
\end{enumerate}

\bitmapfigure[width=0.42\linewidth]{img/1994/national_liberal/q23_fig1.png}
\end{problem}

\begin{solution}
\begin{enumerate}
\item
取底面为 \(z=0\)，令 \(B=(-1,0,0)\)、\(C=(1,0,0)\)、\(A=(0,\sqrt3,0)\)，设棱柱高为 \(h>0\). 于是 \(B_1=(-1,0,h)\)、\(C_1=(1,0,h)\)，且 \(D=\left(\dfrac12,\dfrac{\sqrt3}{2},0\right)\).

在矩形 \(B_1BCC_1\) 中，设两条对角线交点为 \(E\)，则 \(E\) 是 \(BC_1\) 的中点，故 \(E\) 在平面 \(DBC_1\) 内. 又
\(\overrightarrow{DE}=\left(-\dfrac12,-\dfrac{\sqrt3}{2},\dfrac h2\right)\)，而 \(\overrightarrow{AB_1}=\left(-1,-\sqrt3,h\right)=2\overrightarrow{DE}\). 所以 \(AB_1\parallel DE\)，且直线 \(DE\) 在平面 \(DBC_1\) 内，从而 \(AB_1\parallel\) 平面 \(DBC_1\).

 \item
由 \(AB_1\perp BC_1\)，有 \(\overrightarrow{AB_1}\cdot\overrightarrow{BC_1}=(-1,-\sqrt3,h)\cdot(2,0,h)=-2+h^2=0\)，所以 \(h=\sqrt2\). 设 \(F=(0,0,0)\) 为 \(A\) 在 \(BC\) 上的垂足，则 \(F\) 是 \(BC\) 的中点，且 \(AF\perp BC\). 由于侧面 \(B_1BCC_1\) 是平面 \(y=0\)，有 \(AF\perp\) 平面 \(B_1BCC_1\)，故线段 \(AB_1\) 在该侧面上的射影为 \(FB_1\). 于是
\(FB_1=\sqrt{FB^2+BB_1^2}=\sqrt{1^2+(\sqrt2)^2}=\sqrt3\).

所以射影长为 \(\sqrt3\).
\end{enumerate}
\end{solution}

\begin{problem}
已知直角坐标平面上点 \(Q(2,0)\) 和圆 \(C\)：\(x^{2} + y^{2} = 1\)，动点 \(M\) 到圆 \(C\) 的切线长与 \(|MQ|\) 的比等于常数 \(\lambda\)（\(\lambda > 0\)）. 求动点 \(M\) 的轨迹方程，说明它表示什么曲线.
\end{problem}

\begin{solution}
设 \(M=(x,y)\)，圆心为 \(O=(0,0)\)，从 \(M\) 向圆作切线，切点为 \(N\). 由 \(ON\perp MN\) 及 \(ON=1\)，切线长满足 \(MN^2=MO^2-1=x^2+y^2-1\). 题设给出 \(MN=\lambda|MQ|\)，因此
\(x^2+y^2-1=\lambda^2\bigl((x-2)^2+y^2\bigr)\). 整理得轨迹方程
\((\lambda^2-1)(x^2+y^2)-4\lambda^2x+1+4\lambda^2=0\). 反过来，满足此方程的点有 \(x^2+y^2-1=\lambda^2|MQ|^2\geqslant0\)，且不会出现 \(M=Q\)，所以确实都能作相应切线，方程没有引入多余点.

当 \(\lambda=1\) 时，方程化为 \(x=\dfrac54\)，表示一条直线. 当 \(\lambda\neq1\) 时，配方得
\(\left(x-\dfrac{2\lambda^2}{\lambda^2-1}\right)^2+y^2=\dfrac{3\lambda^2+1}{(\lambda^2-1)^2}\)，表示圆，圆心为 \(\left(\dfrac{2\lambda^2}{\lambda^2-1},0\right)\)，半径为 \(\dfrac{\sqrt{3\lambda^2+1}}{|\lambda^2-1|}\).
\end{solution}

\begin{problem}
设数列 \(\{a_n\}\) 的前 \(n\) 项和为 \(S_n\)，若对于所有的自然数 \(n\)，都有 \(S_n = \frac{n(a_1 + a_n)}{2}\)，证明 \(\{a_n\}\) 是等差数列.
\end{problem}

\begin{solution}
令 \(d=a_2-a_1\). 当 \(n=1\) 时，\(a_1=a_1+(1-1)d\) 成立；当 \(n=2\) 时，\(a_2=a_1+d\) 也成立.

假设对某个 \(k\geqslant2\) 有 \(a_k=a_1+(k-1)d\). 由题设及 \(S_{k+1}=S_k+a_{k+1}\)，得
\(\dfrac{(k+1)(a_1+a_{k+1})}{2}=\dfrac{k(a_1+a_k)}{2}+a_{k+1}\). 整理为 \((k-1)a_{k+1}=ka_k-a_1\). 代入归纳假设，得
\((k-1)a_{k+1}=k\bigl[a_1+(k-1)d\bigr]-a_1=(k-1)(a_1+kd)\). 因 \(k-1>0\)，所以 \(a_{k+1}=a_1+kd\).

由数学归纳法，\(a_n=a_1+(n-1)d\) 对所有自然数 \(n\) 成立，因此 \(\{a_n\}\) 是等差数列，公差为 \(d=a_2-a_1\).
\end{solution}
